\section[Jacobian 判别法]{Jacobian 判别法（1 月 15 日）}

% 以保证与原文编号逐一相同（不写死数字）。

\srcpage{102}

设 $E\supset\cO_E$、$\pi$、$\Fq$ 照常。回忆光滑性的 Jacobian 判别法（Jacobian criterion）的陈述：

\begin{theorem}[Jacobian 判别法]
设 $S$ 是 $\Fq$ 上的完满空间。设 $Z\to X_S=X_{S,E}$ 是（sousperfectoid）adic 空间之间的光滑映射，且 $Z$ 拟射影\footnote{即存在 Zariski 闭嵌入 $Z\hookrightarrow U$，其中 $U$ 在 $\PP^n_{X_S}$ 中开。}。设 $\cM_Z$ 是 $Z\to X_S$ 的截面空间，$\cM_Z^{\mathrm{smooth}}\subset\cM_Z$ 是使 $s^*T_{Z/X_S}$ 的 HN 斜率处处为正的开子空间。则映射 $\cM_Z^{\mathrm{smooth}}\to S$ 对一切 $\ell\ne p$ 都是 $\ell$-上同调光滑的。
\end{theorem}

\begin{remark}
回忆 $\cM_Z\to S$ 可表为局部空间菱形、可紧化，且局部有限 dim.tr.g.。这些性质都被开子集 $\cM_Z^{\mathrm{smooth}}$ 继承。
\end{remark}

证明策略如下：
\begin{enumerate}[label=(\arabic*)]
\item 证明 $\cM_Z^{\mathrm{smooth}}\to S$ 的形式光滑性；
\item 证明形式光滑性连同「几何有限维性」蕴含 $\Fl$ 是 $f$-ULA 的；
\item 证明 $Rf^!\Fl$ 可逆。
\end{enumerate}

\subsection{第 (1) 步}

按定义，$\cM_Z^{\mathrm{smooth}}\to S$ 形式光滑这一陈述意味着下述内容。对仿射完满空间的任意 Zariski 闭浸入 $T_0\subset T$，设给定图表
\[
\begin{tikzcd}
T_0 \arrow[r] \arrow[d] & \cM_Z^{\mathrm{smooth}} \arrow[d] \\
T \arrow[r] & S
\end{tikzcd}
\]
则需要证明：存在 étale 映射 $T'\to T$ 使其像包含 $T_0$，并且存在提升
\[
\begin{tikzcd}
T_0\times_T T' \arrow[r] \arrow[d] & \cM_Z^{\mathrm{smooth}} \arrow[d] \\
T' \arrow[r] \arrow[ur, dashed] & S
\end{tikzcd}
\]
证明是对 Banach 不动点定理的一番繁琐的显式分析。

\subsection{第 (2) 步}

若 $\cM_Z^{\mathrm{smooth}}$ 能由一个完满空间表出，则可取 $T_0$ 为 $\cM_Z^{\mathrm{smooth}}$ 的开子集，并借此看出 $\cM_Z^{\mathrm{smooth}}$ 是某个在完满球上 étale 的对象的收缩。

\begin{proposition}
存在上同调光滑且形式光滑的满射 $T_0\to\cM_Z$（可表为局部空间菱形、可紧化，且局部有限 dim.tr.g.），使得 $T_0$ 是完满空间，并且局部容许一个到 $S$ 上有限维完满球的 Zariski 闭嵌入。
\end{proposition}

\begin{remark}
当然，任何菱形都有以完满空间为图册的覆盖，但该图册未必上同调光滑且形式光滑。实践中出现的菱形往往有这样的图册。
\end{remark}

\srcpage{103}

\begin{proof}
映射 $\cM_Z\hookrightarrow\cM_{\PP^n_S}$ 局部 Zariski 闭。这允许把命题化归为射影空间的情形。此时有
\[
\cM_{\PP^n_S}=\bigcup_{d\ge0}\bigl(\mathrm{BC}(\cO(d)^{n+1})\setminus\{0\}\bigr)/E^\times.
\]
于是只需证明 $\mathrm{BC}(\cO(d)^{n+1})/E^\times$ 具有该性质。这可以显式地做到。
\end{proof}

\begin{corollary}
常值层 $\Fl$ 对 $\cM_Z^{\mathrm{smooth}}\to S$ 是 ULA 的。
\end{corollary}

\begin{proof}
ULA 性可在上同调光滑局部化之后检验。故只需对 $T_0\to S$ 检验，而后者由下面的引理推出。
\end{proof}

\begin{lemma}
设 $T_0\to S$ 是仿射完满空间之间的映射，且满足
\begin{enumerate}[label=(\arabic*)]
\item $T_0$ 在 $S$ 上形式光滑；
\item $T_0\hookrightarrow\mathbb{B}^n_S$ 在某个有限维完满球中是 Zariski 闭的。
\end{enumerate}
则 $\Fl$ 对 $T_0\to S$ 是 ULA 的。
\end{lemma}

\begin{proof}
由形式光滑性，可找到 $T:=\mathbb{B}^n_S$ 上的 étale 映射 $T'$，连同图表
\[
\begin{tikzcd}
T_0\times_T T' \arrow[r, "\text{étale}"] \arrow[d] & T_0 \arrow[d] \\
T' \arrow[r, "\text{étale}"'] & T=\mathbb{B}^n_S
\end{tikzcd}
\]
收缩 $T'$，甚至可找到收缩
\[
\begin{tikzcd}
T_0\times_T T' \arrow[r, "\text{étale}"] & T_0 \\
T' \arrow[u, dashed, "r"'] \arrow[d, "\text{étale}"'] & \\
\mathbb{B}^n_S \arrow[r] & S
\end{tikzcd}
\]
于是 $T_0\times_T T'\to S$ 是 $T'\xrightarrow{\text{étale}}\mathbb{B}^n_S\to S$ 的收缩。反过来，$T_0\times_T T'$ 是 $T_0$ 的 étale 覆盖。

我们断言：ULA 性在收缩下保持。（注意，上同调光滑性在收缩下并不保持。）承认这一点，便得到 $T_0\times_T T'\to S$ 是 ULA 的。由于 ULA 性可 étale 局部地检验，这就说明 $T_0\to S$ 也是 ULA 的。

断言的证明：直接由定义即可验出；更巧妙的论证是它由定理 18.30 中 ULA 的范畴刻画推出。
\end{proof}

\subsection{第 (3) 步}

我们要证明 $Rf^!\Fl$ 可逆，即局部同构于 $\Fl[n]$。

\begin{fact}
若 $A$ 对 $f\colon X\to S$ 是 $f$-ULA 的，则 $\mathbb{D}_{X/S}(A)$ 仍是 $f$-ULA 的，并且其构造与任意基变换 $S'\to S$ 交换。
\end{fact}

由于层可逆这一性质可在 $v$-覆盖之后检验，故只需 $v$-局部地检验。由上述事实，对偶在任意基变换下保持。

\srcpage{104}

我们将化归为沿截面拉回之后再证明它。设有图表
\[
\begin{tikzcd}
S' \arrow[r, "g"] \arrow[dr, "h"'] & \cM_Z^{\mathrm{smooth}} \arrow[d, "f"] \\
& S
\end{tikzcd}
\]
其中 $g$ 是 $v$-覆盖。如上所述，只需证明 $g^*Rf^!\Fl$ 可逆。$g$ 的数据等价于一个截面 $s'\colon X_{S'}\to\cM_{Z'}$，其中 $Z':=Z\times_{X_S}X_{S'}$。于是 $g=h's'$，其中 $h'$ 是 $h$ 的基变换，如下图。

\[
\begin{tikzcd}
\cM^{\mathrm{smooth}}_{Z'} \arrow[r, "h'"] \arrow[d, "f'"'] & \cM^{\mathrm{smooth}}_{Z} \arrow[d, "f"] \\
S' \arrow[r, "h"'] \arrow[u, "s'", bend left=30] & S
\end{tikzcd}
\]

则 $g^*Rf^!\Fl\cong(s')^*(h')^*Rf^!\Fl\cong(s')^*R(f')^!\Fl$，最后一步用到上述事实。

改名之后可见：只需对任意截面 $s\colon S\to\cM_Z^{\mathrm{smooth}}$（它可等同于截面 $X_S\to Z$）证明拉回 $s^*Rf^!\Fl\in D_{\et}(S,\Fl)$ 可逆。

现在沿此截面实现法锥的形变。
\[
\begin{tikzcd}
Z \\
X_S \arrow[u, "s"']
\end{tikzcd}
\]

由此可以造出
\[
\begin{tikzcd}
\wt Z \arrow[d, "\text{smooth}"] \\
X_S\times\mathbf{A}^1
\end{tikzcd}
\]
（其上带截面 $\wt s$）使得
\begin{itemize}
\item $\wt Z\times_{\mathbf{A}^1}\{1\}=Z$；
\item $\wt Z\times_{\mathbf{A}^1}\{0\}$ 是 $s$ 在 $Z$ 中的法锥，即对应于 $s^*T_{Z/X_S}$ 的几何向量丛。
\end{itemize}
我们想把整族实现为同一问题的实例，然而基 $X_S\times\mathbf{A}^1$ 不再是我们先前考虑的 $X_{\ul S}$ 的形状。在 $X_S\times\mathbf{A}^1$ 内有 $X_S\times\ul E=X_{S\times\ul E}$，它带 $\ul E^\times$-作用。拉回 $\wt Z\to X_S\times\mathbf{A}^1$ 得到 $\wt Z'\to X_{S\times\ul E}$。这是光滑的，它在 $S\times\{1\}$ 上的纤维是 $Z$，在 $S\times\{0\}$ 上的纤维是 $s^*T_{Z/X_S}$。

$\wt Z'$ 仍是拟射影的，故先前的全部结果都适用。于是有
\[
\begin{tikzcd}
\cM^{\mathrm{smooth}}_{\wt Z'} \arrow[r, "\wt f", bend left=20] & S\times\ul E \arrow[l, "\wt s", bend left=20]
\end{tikzcd}
\]

则 $R\wt f^!\Fl$ 是 $\wt f$-ULA 的（它是 $\wt f$-ULA 层 $\Fl$ 的 Verdier 对偶），并且
\begin{itemize}
\item $R\wt f^!\Fl|_{S\times\{1\}}\cong Rf^!\Fl$；
\item $R\wt f^!\Fl|_{S\times\{0\}}$ 是 $\mathrm{BC}(s^*T_{Z/X_S})$ 的对偶复形，而由显式计算知它可逆。
\end{itemize}

\srcpage{105}

由一番巧妙的展布论证可知，$s^*Rf^!\Fl$ 也可逆。（我们知道形变上的整个对偶复形是 ULA 的。）

\subsection{对 $D_{\et}(\Bun_G,\Lambda)$ 的应用}

回忆我们按如下方式为 $\Bun_G$ 定义图卡。

\begin{definition}
设 $\cM$ 是 $\mathbf{Q}$-滤过 $G$-丛（滤过递增）的模空间，即正合 $\otimes$-函子
\[
\Rep_E(G)\xrightarrow{\ \rho\ }\mathbf{Q}\,\mathrm{Fil}_{\Bun_{X_S}},
\]
使得对一切 $V\in\Rep_E(G)$，$\rho(V)^{\le\lambda}\big/\bigcup_{\lambda'<\lambda}\rho(V)^{\le\lambda'}$ 是斜率为 $\lambda$ 的半稳定对象。（这些就是「反向 HN 滤过」。）
\end{definition}

存在一个遗忘滤过的映射 $\cM\to\Bun_G$。另一方面，把滤过送到其伴随分次给出映射 $\cM\to\bigcup_{b\in B(G)}[\ast/G_b(E)]$，因为半稳定丛等价于等晶体。这诱导出分解
\[
\cM=\bigcup_{b\in B(G)}\cM_b.
\]

\begin{theorem}
映射 $\cM\to\Bun_G$ 上同调光滑。
\end{theorem}

\begin{example}
设 $G=\mathrm{GL}_2$，$b\leftrightarrow\cO\oplus\cO(1)$。则 $\cM_b$ 参数化扩张
\[
0\to\cL\to\cE\to\cL'\to0,
\]
其中 $\deg\cL=0$，$\deg\cL'=1$。
\end{example}

定理 19.9 是 Jacobian 判别法的推论。取 $S\to\Bun_G$ 对应于 $G$-丛 $\cE/X_S$，并取 $Z$ 为 $\cE$ 上 $\mathbf{Q}$-滤过的模空间。则 $\cM\hookrightarrow\cM_Z$，而由斜率的条件它实际上落在 $\cM_Z^{\mathrm{smooth}}$ 中。

\subsection{$\cM_b$ 的几何}

现在固定 $b\in B(G)$，考虑 $\pi_b\colon\cM_b\to\Bun_G$。我们把它看作「$\Bun_G$ 在 $\Bun^b_G$ 附近的图卡」。

由于有映射 $\cM_b\to[\ast/G_b(E)]$，可把 $\cM_b$ 看作 $[\wt\cM_b/G_b(E)]$，其中在 $\wt\cM_b$ 里分次丛被平凡化。

\begin{example}
设 $G=\mathrm{GL}_2$，$b\leftrightarrow\cO\oplus\cO(1)$。则 $\wt\cM_b$ 参数化扩张
\[
0\to\cO\to\cE\to\cO(1)\to0.
\]
\end{example}

存在一个「基点」$\ast\in\wt\cM_b$，对应于分裂扩张。这给出截面
\[
\begin{tikzcd}
{}[\ast/G_b(E)] \arrow[dr, "\sim"'] & \cM_b \arrow[l] \arrow[d] \\
& {}[\ast/G_b(E)]
\end{tikzcd}
\]

事实上，图表
\[
\begin{tikzcd}
{}[\ast/G_b(E)] \arrow[r] \arrow[d] & \cM_b \arrow[d] \\
\Bun^b_G \arrow[r] & \Bun_G
\end{tikzcd}
\]
是笛卡的。因为若 $\cE\in\Bun^b_G$，则 $\cE$ 的 HN 滤过给出给定 $\mathbf{Q}$-滤过的分裂。

\srcpage{106}
\begin{center}
\figP{106}
\end{center}

\begin{fact}
映射 $\wt\cM_b\to\ast$ 可表为局部空间菱形，并且上同调光滑。更确切地说，它是负 Banach–Colmez 空间的逐次扩张，总维数为 $\langle2\rho,\nu_b\rangle$。
\end{fact}

\begin{example}
对如例 19.11 的 $b$，$\cM_b=\{0\to\cO\to\cE\to\cO(1)\to0\}$ 可等同于 $\mathrm{BC}(\cO(-1)[1])$。
\end{example}

\begin{fact}
$\wt\cM_b\setminus\{\ast\}=:\wt\cM_b^\circ$ 是空间菱形。（特别地它是 qcqs 的，但不在 $\ast$ 上 qcqs。）
\end{fact}

\begin{example}
对如例 19.11 的 $b$，有 $\wt\cM_b^\circ\cong\Spa\ol{\Fq}((t^{1/p^\infty}))/\mathrm{SL}_1(D)$，其中 $D/E$ 是四元数代数。这是仿射完满空间对射有限群的商。

事实上，在 $\wt\cM_b^\circ$ 上有 $\cE\cong\cO(1/2)$。取定这样一个同构，则 $\cO\hookrightarrow\cO(1/2)$ 给出 $\mathrm{BC}(\cO(1/2))\setminus\{0\}=\Spa\ol{\Fq}((t^{1/p^\infty}))$ 的一个截面。这类 torsor（与分次丛的选定平凡化相容）构成 $\mathrm{SL}_1(D)$-torsor，故 $\wt\cM_b^\circ$ 是商。
\end{example}

\begin{warning}
在上例中，被挖去的点 $\ast$ 位于 $|t|=1$ 处，并不在 $|t|=0$ 附近。注意当 $t\to0$ 时，映射趋于 $0$，而不是趋于该扩张。
\end{warning}

$\wt\cM_b$ 在下述意义下是「严格局部」的。

\begin{theorem}
对任意 $A\in D_{\et}(\wt\cM_b,\Lambda)$，限制
\[
R\Gamma(\wt\cM_b,A)\to R\Gamma(\ast,A)
\]
是同构。
\end{theorem}

\begin{proof}[梗概]
该映射的锥是「部分紧支撑上同调」$R\Gamma_{\partial c}(\wt\cM_b^\circ,A)$，其中 $R\Gamma_{\partial c}$ 表示「朝 $\ast$ 紧支撑，但对 $\wt\cM_b$ 的边界不加支撑条件」。这是一个远为一般的现象的特例：

设 $X$（我们心里想的是 $X=\wt\cM_b^\circ$）是 $\ol{\Fq}$ 上 dim.tr.g. $<\infty$ 的空间菱形。假设映射 $X\to\ast$ 部分紧合，即 $X(R,R^+)=X(R,R^\circ)$。则对任意 $C/\Fq$，$X_C$ 有「两个端点」。
\end{proof}

\begin{example}
设 $X=\Spa(R,R^+)$ 是仿射完满的，$C=\ol{\Fq}((t))^{\wedge}$。则 $X_C$ 是 $X\times_{\Fq}\Spa\Fq((t))$ 的射有限覆盖，而后者是 $X$ 上的去心开单位圆盘。它有两个端点（原点与外边界）。
\end{example}

\srcpage{107}
\begin{center}
\figP{107}
\end{center}

于是可以定义「部分紧支撑上同调」$R\Gamma_{\partial c}(X_C,A)$：它在一个端点处紧支撑，在另一个端点处则不然。

\begin{theorem}
在上述情形，$R\Gamma_{\partial c}(X_C,A)=0$。
\end{theorem}

我们将给出两条论证：一条纯哲学，一条沿实际证明的思路。实际证明：利用「对应」（即把一般情形前推到这一情形）与紧合基变换，化归到
\[
X=\Spa\ol{\Fq}((t^{1/p^\infty}))
\]
的情形。由于该空间很简单，可以再化归到 $A=\Lambda$ 并直接计算。

直观图像如下：设 $M$ 是带 $\mathbf{R}$ 的自由作用（一个「流」）的拓扑流形，且 $\ol M=M/\mathbf{R}$ 紧。

有两个边界：流的「源」与流的「汇」。对一切 $A\in D(M/\mathbf{R},\mathbf{Z})$，有 $R\Gamma_{\partial c}(M,A)=0$。的确，粘入一个边界圆盘之后，流便收缩到边界。

这如何类比？粗略地说，令 $C=\ol{\Fq}((t^{\mathbf{R}}))$。它有 $\mathbf{R}_{>0}$ 通过重标度的作用（尽管这不连续）。于是 $X_C$ 带 $\mathbf{R}$-作用，扮演 $M$ 的角色；而商 $X_C/\mathbf{R}$ 是 qcqs 的。
